\section{Récupération de l’état et de sa dynamique}
\label{nuclear:3}
\subsection{Représentation récupérable du transport}

Soit \(C\) la réalisation unitaire du transport de mémoire sur \(\mathcal H\). Le système de coordonnées à neuf phases, son inclusion uniforme et sa compression sont

\[
S_C(v_0,\ldots,v_8)=(Cv_8,v_0,\ldots,v_7),\qquad
Jv=(v,\ldots,v)/3,\qquad T=(8I+C)/9.
\]

Soit \(\eta=(I-JJ^*)S_CJ\). Son action est

\[
\eta v=\frac1{27}\bigl(8(C-I)v,-(C-I)v,\ldots,-(C-I)v\bigr).
\]

La normalisation de neuf copies détermine \(1/3\) ; la répartition de la couture détermine \(8/81\) dans

\begin{equation}
I-T^*T=\eta^*\eta=\frac8{81}(I-C)^*(I-C).
\label{nuc03:eq:1}
\end{equation}

Nous écrivons \(P=\operatorname{proj}\ker(I-C)\). Le théorème précédent prouve la convergence forte \(T^N\to P\). Par conséquent,

\begin{equation}
\mathcal Vv=(Pv,\eta v,\eta Tv,\eta T^2v,\ldots)
\quad\text{en}\quad
\mathcal R=P\mathcal H\oplus\ell^2(\mathbb N_0;\mathcal H^9)
\label{nuc03:eq:2}
\end{equation}

est une isométrie. Nous notons son image fermée \(\mathcal M\).

Dans l’espace des registres, nous définissons le décalage de blocs


\begin{equation}
\mathcal B(a,m_0,m_1,m_2,\ldots)=(a,m_1,m_2,\ldots).
\label{nuc03:eq:3}
\end{equation}

La première composante conserve le secteur fixe ; les autres composantes sont les registres vectoriels complets, y compris leurs corrélations.

\begin{theorem}
\label{nuc03:transporte}
Le transport d’origine a la représentation exacte

\begin{equation}
\boxed{\mathcal V C=(9\mathcal B-8I)\mathcal V.}
\label{nuc03:eq:4}
\end{equation}

La restriction \(\mathcal C=(9\mathcal B-8I)|_{\mathcal M}\) est unitaire sur \(\mathcal M\), et


\begin{equation}
C=\mathcal V^*\mathcal C\mathcal V.
\label{nuc03:eq:5}
\end{equation}
\end{theorem}

\begin{proof}
Comme \(PT=P\), la définition~\eqref{nuc03:eq:2} donne composante par composante

\[
\mathcal VT v=(Pv,\eta Tv,\eta T^2v,\ldots)=\mathcal B\mathcal Vv.
\]

L’identité algébrique \(C=9T-8I\) implique \eqref{nuc03:eq:4}. L’isométrie \(\mathcal V:\mathcal H\to\mathcal M\) est surjective par définition de son image. L’égalité \eqref{nuc03:eq:4} identifie la restriction à \(\mathcal V C\mathcal V^{-1}\), qui est unitaire parce que \(C\) l’est. Cela démontre en outre que \(\mathcal C\mathcal M=\mathcal M\) et \eqref{nuc03:eq:5}.

\end{proof}

La restriction au sous-espace des registres compatibles est essentielle : \(9\mathcal B-8I\) sur tout \(\mathcal R\) n’est pas unitaire. Le théorème caractérise le transport sur les registres qui procèdent effectivement d’une même trajectoire linéaire du système de coordonnées documenté.

\begin{corollary}
Pour \(v,w\in\mathcal H\) et \(k\in\mathbb Z\),

\begin{equation}
\langle C^kv,w\rangle
=\langle\mathcal C^k\mathcal Vv,\mathcal Vw\rangle.
\label{nuc03:eq:6}
\end{equation}

Tous les moments croisés de la dynamique réalisée sont ainsi conservés. Les puissances négatives appartiennent à l’inverse de \(\mathcal C\) dans \(\mathcal M\), et non à une inversion du décalage unilatéral sur l’espace ambiant.
\end{corollary}

\begin{corollary}
Pour tout observable borné \(A\), sa réalisation dans l’image est \(\widehat A=\mathcal V A\mathcal V^*|_{\mathcal M}\). Les produits, les adjoints et les commutateurs sont préservés. En particulier,

\begin{equation}
[A,C]=0\quad\Longleftrightarrow\quad[\widehat A,\mathcal C]=0.
\label{nuc03:eq:7}
\end{equation}
\end{corollary}

La preuve consiste à utiliser \(\mathcal V^*\mathcal V=I\) et le fait que \(\mathcal V\mathcal V^*\) est l’identité de \(\mathcal M\). La conservation de la symétrie s’exprime par le transport intégral de ses opérateurs, en plus d’une égalité de normes.


\subsection{Reconstruction explicite et profondeur du registre}

Soit \(m_j=\eta T^jv\). L’identité télescopique donne pour chaque \(N\)

\begin{equation}
\boxed{v=T^{*N}T^Nv+\sum_{j=0}^{N-1}T^{*j}\eta^*m_j.}
\label{nuc03:eq:8}
\end{equation}

En passant à la limite forte,

\begin{equation}
\boxed{v=Pv+\sum_{j=0}^{\infty}T^{*j}\eta^*m_j.}
\label{nuc03:eq:9}
\end{equation}

La reconstruction partielle avec terminal fixe,
\(v_N=Pv+\sum_{j<N}T^{*j}\eta^*m_j\), a pour erreur exacte

\begin{equation}
v-v_N=T^{*N}T^N(I-P)v.
\label{nuc03:eq:10}
\end{equation}

La série converge pour chaque état. Si le registre complet présente une erreur \(\delta r\) en norme de \(\mathcal R\), sa reconstruction au moyen de \(\mathcal V^*\) varie d’une norme au plus égale à \(\|\delta r\|\), car \(\|\mathcal V^*\|=1\). C’est une stabilité uniforme de la lecture complète.

\subsubsection{Différence entre profondeur finie et profondeur illimitée}

Pour le décalage bilatéral \(Ce_k=e_{k+1}\) dans \(\ell^2(\mathbb Z)\), on a \(P=0\). Considérons


\[
v_L=L^{-1/2}\sum_{k=1}^{L}e_k.
\]

Alors \(\|(I-C)v_L\|^2=2/L\). La commutation de \(T\) avec \(I-C\), la contraction de \(T\) et \eqref{nuc03:eq:1} impliquent

\begin{equation}
\sum_{j=0}^{N-1}\|\eta T^jv_L\|^2
\leq\frac{16N}{81L}.
\label{nuc03:eq:11}
\end{equation}

Pour toute profondeur fixe \(N\), le membre de droite tend vers zéro lorsque \(L\to\infty\), bien que \(\|v_L\|=1\). Le registre de mémoire de profondeur fixe, séparé du terminal, ne possède pas de borne inférieure uniforme. En revanche, le registre illimité a une norme exactement égale à un pour chaque \(v_L\).


La différence correspond à deux quantificateurs distincts : chaque état est conservé lorsqu’on complète son registre ; une profondeur fixée à l’avance ne retrouve pas uniformément tous les états sans leur terminal. La conservation à une profondeur arbitraire possède ici une expression fonctionnelle exacte.

\subsubsection{Information vectorielle et intensité}

Le théorème conserve \(m_j\), et non uniquement \(\|m_j\|^2\). Par exemple, dans la réalisation bilatérale, \(v=e_0\) et \(w=e_1\) ont des intensités identiques à chaque profondeur parce que le décalage commute avec \(T\) et transporte \(\eta\) unitairement. Leurs registres vectoriels sont différents et \eqref{nuc03:eq:9} récupère les deux états distincts. Les intensités servent au bilan ; les registres avec leurs corrélations servent en outre à la reconstruction.

\subsection{Naturalité de la récupération sous raffinement
}

Soient \(\iota_n:\mathcal H_n\to\mathcal H_{n+1}\) les isométries de raffinement obtenues à partir des émissions communes, et \(C_n\) des réalisations unitaires avec

\begin{equation}
C_{n+1}\iota_n=\iota_nC_n.
\label{nuc03:eq:12}
\end{equation}

La relation vaut également pour les adjoints : multiplier par \(C_{n+1}^*\) et \(C_n^*\) donne \(C_{n+1}^*\iota_n=\iota_nC_n^*\). Par conséquent, les images sont réductrices. Les formules polynomiales de \(T_n\), \(\eta_n\) et la limite forte de \(T_n^N\) donnent


\begin{equation}
T_{n+1}\iota_n=\iota_nT_n,\quad
\eta_{n+1}\iota_n=\iota_n^{\oplus9}\eta_n,\quad
P_{n+1}\iota_n=\iota_nP_n.
\label{nuc03:eq:13}
\end{equation}

Soit \(\mathfrak I_n\) le raffinement qui applique \(\iota_n\) au terminal et \(\iota_n^{\oplus9}\) à chaque registre. Alors

\begin{equation}
\boxed{\mathcal V_{n+1}\iota_n=\mathfrak I_n\mathcal V_n,\qquad
\mathcal B_{n+1}\mathfrak I_n=\mathfrak I_n\mathcal B_n.}
\label{nuc03:eq:14}
\end{equation}

\begin{proof}
Chaque composante de la première égalité est une identité de \eqref{nuc03:eq:13}, itérée à la puissance correspondante de \(T_n\). La seconde égalité affirme que le raffinement de chaque bloc commute avec le décalage de l’indice des blocs, ce qui se vérifie directement. Les deux applications sont bornées, de sorte que les égalités passent des registres à support fini à leur complétion.
\end{proof}

Par \eqref{nuc03:eq:4} et \eqref{nuc03:eq:14}, la reconstruction de la dynamique commute aussi avec le raffinement. L’indice de profondeur généalogique \(n\) et celui des compressions successives \(j\) restent distincts. L’égalité exprime leur compatibilité lorsqu’ils agissent au moyen des carrés documentés, au lieu de les identifier nominalement.

\subsection{Deux régimes exacts de distribution de mémoire
}

Pour un état unitaire \(v\) tel que \(Pv=0\), nous définissons


\begin{equation}
a_N=\|T^Nv\|^2,\qquad b_N=\|\eta T^Nv\|^2=a_N-a_{N+1}.
\label{nuc03:eq:15}
\end{equation}

On a \(a_0=1\), \(a_N\to0\), \(b_N\geq0\) et \(\sum b_N=1\). Les poids \(b_N\) décrivent la distribution entre profondeurs du registre de mémoire. Cet indice compte des compressions ; son identification à une durée physique requerrait d’appliquer le lecteur temporel correspondant.


\subsubsection{Transport exceptionnel d’ordre trois}

L’opérateur orthogonal sans point fixe \(g_c\) satisfait \(g_c^2+g_c+I=0\). Pour la composition explicite \(C=g_c\), le développement précédent a obtenu

\begin{equation}
T^*T=\frac{19}{27}I,
\qquad a_N=\left(\frac{19}{27}\right)^N,
\qquad b_N=\frac8{27}\left(\frac{19}{27}\right)^N.
\label{nuc03:eq:16}
\end{equation}

La profondeur moyenne, en comptant le premier registre comme profondeur un, est

\begin{equation}
\sum_{N\geq0}(N+1)b_N=\frac{27}{8}.
\label{nuc03:eq:17}
\end{equation}

C’est un résultat de cette composition du système de coordonnées nonadique avec le transport exceptionnel déjà construit ; il conserve le domaine et n’identifie pas, par son ordre, \(g_c\) à l’holonomie temporelle complète.

\subsubsection{Registre bilatéral de tours}

Pour \(C e_k=e_{k+1}\) et \(v=e_0\), le binôme détermine tous les termes :

\begin{equation}
T^Ne_0=9^{-N}\sum_{k=0}^{N}\binom{N}{k}8^{N-k}e_k,
\quad
a_N=81^{-N}\sum_{k=0}^{N}\binom{N}{k}^{\!2}64^{N-k}.
\label{nuc03:eq:18}
\end{equation}

Cette formule est exacte à toute profondeur. La représentation circulaire ultérieure fournit

\begin{equation}
a_N=\frac1{2\pi}\int_{-\pi}^{\pi}
\left(\frac{65+16\cos\theta}{81}\right)^N\,d\theta.
\label{nuc03:eq:19}
\end{equation}

\begin{theorem}
Dans cette réalisation bilatérale,

\begin{equation}
\boxed{a_N\sim\frac9{4\sqrt{2\pi}}N^{-1/2},\qquad
b_N\sim\frac9{8\sqrt{2\pi}}N^{-3/2}.}
\label{nuc03:eq:20}
\end{equation}
\end{theorem}

\begin{proof}
Nous écrivons \(r(\theta)=(65+16\cos\theta)/81\) et \(q=8/81\). Au voisinage de zéro,

\[
r(\theta)=1-q\theta^2+O(\theta^4),\qquad
\log r(\theta)=-q\theta^2+O(\theta^4).
\]

Un voisinage suffisamment petit étant fixé, il existe des constantes positives \(c_1,c_2\) qui encadrent \(-\log r(\theta)\) entre \(c_1\theta^2\) et \(c_2\theta^2\). Hors de ce voisinage, \(r\leq\rho<1\). Le changement \(x=\sqrt N\theta\), la borne gaussienne et la convergence dominée donnent

\[
\lim_{N\to\infty}\sqrt N\,a_N
=\frac1{2\pi}\int_{\mathbb R}e^{-qx^2}\,dx
=\frac1{2\sqrt{\pi q}}=\frac9{4\sqrt{2\pi}}.
\]

Pour la seconde limite, on applique directement le même procédé à


\[
b_N=\frac1{2\pi}\int_{-\pi}^{\pi}r(\theta)^N(1-r(\theta))\,d\theta.
\]

Dans la variable \(x\), \(N(1-r(x/\sqrt N))\to qx^2\), dominé localement par une constante multipliée par \(x^2\) ; hors du voisinage, la contribution est exponentiellement petite même après multiplication par \(N^{3/2}\). Par conséquent

\[
\lim_{N\to\infty}N^{3/2}b_N
=\frac q{2\pi}\int_{\mathbb R}x^2e^{-qx^2}\,dx
=\frac1{4\sqrt{\pi q}}=\frac9{8\sqrt{2\pi}}.
\]

Cette preuve de la seconde limite utilise sa propre intégrale, sans différencier un équivalent asymptotique non contrôlé.

\end{proof}

Par Tonelli et l’identité des queues d’une suite positive,


\begin{equation}
\sum_{N\geq0}(N+1)b_N=\sum_{N\geq0}a_N=+\infty.
\label{nuc03:eq:21}
\end{equation}

Toute la norme est conservée en mémoire, mais sa profondeur moyenne diverge dans cet état bilatéral. La loi de répartition dépend de la représentation effective du transport : le secteur exceptionnel d’ordre trois possède un taux géométrique, tandis que le registre bilatéral considéré ici possède un taux algébrique.

\subsection{Composition avec l’identité d’Euler et l’incidence}

L’identité d’Euler étendue de la section~\ref{euler-trit-generadores} réalise les trois régimes au moyen du Coxeter de \(A_2\), d’un transport nilpotent et de \(F_{\rm av}^2\), avec leurs antécédents déclarés. Leurs générateurs normalisés satisfont \(J_\tau^2=-\tau I\) et ont une trace nulle. Dans leurs systèmes de coordonnées bidimensionnels,

\[
C_\tau(t)=\exp(tJ_\tau),\qquad
\Omega=\begin{pmatrix}0&1\\-1&0\end{pmatrix},\qquad
C_\tau(t)^\mathsf T\Omega C_\tau(t)=\Omega.
\]

La trace nulle donne \(\det C_\tau(t)=1\), et l’identité \(M^\mathsf T\Omega M=(\det M)\Omega\) en dimension deux prouve la dernière égalité. Pour

\[
T_\tau=\frac{8I+C_\tau}{9},\qquad
D_\tau=\frac{\sqrt8}{9}(C_\tau-I),
\]

on obtient une loi de conservation orientée commune :

\begin{equation}
\boxed{\Omega=T_\tau^\mathsf T\Omega T_\tau+D_\tau^\mathsf T\Omega D_\tau.}
\label{nuc03:eq:22}
\end{equation}

\begin{proof}
En développant les deux termes, les contributions croisées de \(C_\tau\) ont des coefficients opposés. Il reste \(72\Omega/81+9C_\tau^\mathsf T\Omega C_\tau/81=\Omega\).
\end{proof}

Cette forme alternée correspond à l’aire orientée du plan canonique. Sa conservation permet de comparer les trois régimes en respectant leurs différences. Pour les représentants concrets \(C\) de Coxeter, \(I+N\) et \(F_{\rm av}^2\), respectivement,

\begin{equation}
(\det T,\det D)=
\left(\frac{19}{27},\frac8{27}\right),\qquad
(1,0),\qquad
\left(\frac{89}{81},-\frac8{81}\right).
\label{nuc03:eq:23}
\end{equation}

La somme vaut un dans les trois cas. Dans le cas parabolique, \(D\) est de rang un et peut conserver des informations même si son aire est nulle. Dans le cas hyperbolique, le complément a une orientation opposée ; le bilan porte sur l’aire signée. Le théorème~\ref{nuc03:transporte} utilise une réalisation unitaire et positive : les hypothèses de chaque résultat sont expressément conservées.


Le développement de l’action et de la forme variable prouve la version entre systèmes de coordonnées différents, son intégration sur des courbes fermées et sa compatibilité avec le raffinement. Le développement de l’incidence compose la reconstruction avec \(F(k)=(\mu(k),\mathcal I P_0k/6)\), en maintenant les douze coordonnées de \(K\) ; il transporte également les opérateurs et leurs repères locaux. Ainsi, une même identité de conservation admet des réalisations typées de norme, d’incidence et d’action orientée.
